\documentclass{article}

\usepackage{amsmath, amssymb, amsthm}
\usepackage[a4paper, margin=1in]{geometry}

\title{Wilson's Theorem}
\author{sid}
\date{\today}

\newtheorem{theorem}{Theorem}
\newtheorem{lemma}{Lemma}
\newtheorem{example}{Example}

\newcommand{\Z}{\mathbb{Z}}
\newcommand{\bld}[1]{\textbf{#1}}
\newcommand{\itl}[1]{\textit{#1}}
\newcommand{\N}{\mathbb{N}}
\newcommand{\Q}{\mathbb{Q}}
\newcommand{\R}{\mathbb{R}}
\newcommand{\eps}{\varepsilon}
\newcommand{\poly}{\mathrm{poly}}
\newcommand{\polylog}{\mathrm{polylog}}
\newcommand{\bigO}{\mathcal{O}}

\begin{document}

\maketitle

\paragraph{Claim.} A natural number $n$ is prime iff
$$
(n-1)! \equiv -1\bmod n
$$

\paragraph{Proof.}
Suppose that $n$ is composite. Then there exists a prime $q$ such that
\[
2 \leq q < n \quad \text{and} \quad q \mid n.
\]
Hence, there exists an integer $k$ such that
\[
n = qk.
\]

Assume for contradiction that
\[
(n-1)! \equiv -1 \pmod{n}.
\]
Then there exists an integer $m$ such that
\[
(n-1)! = nm - 1 = (qk)m - 1 = q(km) - 1.
\]
Thus, $(n-1)!$ is one less than a multiple of $q$, which implies
\[
(n-1)! \equiv -1 \pmod{q}.
\]

On the other hand, since $2 \leq q \leq n-1$, the integer $q$ appears among the factors in
\[
(n-1)! = (n-1)(n-2)\cdots 2 \cdot 1.
\]
Therefore,
\[
(n-1)! \equiv 0 \pmod{q}.
\]

This is a contradiction. Hence,
$
(n-1)! \not\equiv -1 \pmod{n}
$
when $n$ is composite.
\hfill $\square$

\end{document}